%%% Supplementary material of the AAMAS 2027 submission (anonymous). Build: paper/build.sh
\documentclass[10pt]{article}
\usepackage[margin=2.2cm]{geometry}
\usepackage{libertine}
\usepackage[libertine]{newtxmath}
\usepackage[T1]{fontenc}
\usepackage{booktabs}
\usepackage{longtable}
\usepackage{amsmath}
\usepackage{graphicx}
\usepackage{xspace}
\usepackage[numbers,sort&compress]{natbib}
\usepackage{xr-hyper}
\usepackage[hidelinks]{hyperref}
\externaldocument[M-]{main}
\pdfmapfile{+libertine.map}\pdfmapfile{+zi4.map}\pdfmapfile{+newtx.map}
% generated from docs/results/val-c1 and runs/anytime_c1 by paper/tools/make_tables.py
\newcommand{\ConeBcpsatlo}{0.885}
\newcommand{\ConeBcpsathi}{0.968}
\newcommand{\ConeBpcpsatlo}{0.889}
\newcommand{\ConeBpcpsathi}{0.951}
\newcommand{\ConeBcpfulllo}{0.867}
\newcommand{\ConeBcpfullhi}{0.930}
\newcommand{\ConeBctaslo}{0.127}
\newcommand{\ConeBctashi}{0.538}
\newcommand{\ConeBconstructlo}{0.885}
\newcommand{\ConeBconstructhi}{0.934}
\newcommand{\ConeBrllo}{0.680}
\newcommand{\ConeBrlhi}{0.799}
\newcommand{\ConeBclasslo}{3.2}
\newcommand{\ConeBclasshi}{13.3}
\newcommand{\ConeBrlgainlo}{20.1}
\newcommand{\ConeBrlgainhi}{32.0}
\newcommand{\ConeBheld}{8}
\newcommand{\ConeBverdict}{is met}
\newcommand{\ConeBsettings}{8}
\newcommand{\ConeBheldlist}{single-skill binary 25/5/50; single-skill binary 50/5/50; single-skill additive 50/5/50; multi-skill additive 25/5/50; multi-skill additive 50/5/50; multi-skill additive 50/5/200; multi-skill additive 150/10/500; multi-skill additive 150/5/500}
\newcommand{\ConeBfaillist}{none}
\newcommand{\ConeBtests}{46}
\newcommand{\ConeBsig}{46}
\newcommand{\ConeBmaxp}{0.013}
\newcommand{\ConeBwon}{410}
\newcommand{\ConeBpairs}{410}
\newcommand{\ConeBBcpsatlo}{--}
\newcommand{\ConeBBcpsathi}{--}
\newcommand{\ConeBBpcpsatlo}{0.907}
\newcommand{\ConeBBpcpsathi}{0.956}
\newcommand{\ConeBBcpfulllo}{0.886}
\newcommand{\ConeBBcpfullhi}{0.944}
\newcommand{\ConeBBctaslo}{--}
\newcommand{\ConeBBctashi}{--}
\newcommand{\ConeBBconstructlo}{0.899}
\newcommand{\ConeBBconstructhi}{0.943}
\newcommand{\ConeBBrllo}{--}
\newcommand{\ConeBBrlhi}{--}
\newcommand{\ConeBBclasslo}{4.4}
\newcommand{\ConeBBclasshi}{11.4}
\newcommand{\ConeBBrlgainlo}{--}
\newcommand{\ConeBBrlgainhi}{--}
\newcommand{\ConeBBheld}{6}
\newcommand{\ConeBBverdict}{is met}
\newcommand{\ConeBBsettings}{6}
\newcommand{\ConeBBheldlist}{single-skill binary 25/5/50; single-skill binary 50/5/50; single-skill additive 50/5/50; multi-skill additive 25/5/50; multi-skill additive 50/5/50; multi-skill additive 50/5/200}
\newcommand{\ConeBBfaillist}{none}
\newcommand{\ConeBBtests}{17}
\newcommand{\ConeBBsig}{17}
\newcommand{\ConeBBmaxp}{9.4e-05}
\newcommand{\ConeBBwon}{170}
\newcommand{\ConeBBpairs}{170}
\newcommand{\Ctwocpsatlo}{0.934}
\newcommand{\Ctwocpsathi}{0.995}
\newcommand{\Ctwopcpsatlo}{0.927}
\newcommand{\Ctwopcpsathi}{0.987}
\newcommand{\Ctwocpfulllo}{0.920}
\newcommand{\Ctwocpfullhi}{0.982}
\newcommand{\Ctwoctaslo}{0.133}
\newcommand{\Ctwoctashi}{0.553}
\newcommand{\Ctwoconstructlo}{0.933}
\newcommand{\Ctwoconstructhi}{0.995}
\newcommand{\Ctworllo}{0.753}
\newcommand{\Ctworlhi}{0.833}
\newcommand{\Ctwoclasslo}{0.5}
\newcommand{\Ctwoclasshi}{8.0}
\newcommand{\Ctworlgainlo}{16.7}
\newcommand{\Ctworlgainhi}{24.7}
\newcommand{\Ctwoheld}{4}
\newcommand{\Ctwoverdict}{is not met}
\newcommand{\Ctwosettings}{8}
\newcommand{\Ctwoheldlist}{single-skill binary 25/5/50; multi-skill additive 25/5/50; multi-skill additive 50/5/50; multi-skill additive 50/5/200}
\newcommand{\Ctwofaillist}{single-skill binary 50/5/50 (against CP-LNS, Parallel CP-LNS); single-skill additive 50/5/50 (against CP-LNS, Parallel CP-LNS, CP-SAT full, Greedy restarts); multi-skill additive 150/10/500 (against Greedy restarts); multi-skill additive 150/5/500 (against CP-LNS, Parallel CP-LNS)}
\newcommand{\Ctwotests}{46}
\newcommand{\Ctwosig}{37}
\newcommand{\Ctwomaxp}{0.27}
\newcommand{\Ctwowon}{395}
\newcommand{\Ctwopairs}{410}
\newcommand{\CtwoOnecpsatlo}{0.937}
\newcommand{\CtwoOnecpsathi}{1.014}
\newcommand{\CtwoOnepcpsatlo}{0.934}
\newcommand{\CtwoOnepcpsathi}{1.006}
\newcommand{\CtwoOnecpfulllo}{0.923}
\newcommand{\CtwoOnecpfullhi}{1.002}
\newcommand{\CtwoOnectaslo}{0.133}
\newcommand{\CtwoOnectashi}{0.564}
\newcommand{\CtwoOneconstructlo}{0.937}
\newcommand{\CtwoOneconstructhi}{1.014}
\newcommand{\CtwoOnerllo}{0.759}
\newcommand{\CtwoOnerlhi}{0.841}
\newcommand{\CtwoOneclasslo}{-1.4}
\newcommand{\CtwoOneclasshi}{7.7}
\newcommand{\CtwoOnerlgainlo}{15.9}
\newcommand{\CtwoOnerlgainhi}{24.1}
\newcommand{\CtwoOneheld}{3}
\newcommand{\CtwoOneverdict}{is not met}
\newcommand{\CtwoOnesettings}{8}
\newcommand{\CtwoOneheldlist}{multi-skill additive 25/5/50; multi-skill additive 50/5/50; multi-skill additive 50/5/200}
\newcommand{\CtwoOnefaillist}{single-skill binary 25/5/50 (against Parallel CP-LNS); single-skill binary 50/5/50 (against CP-LNS, Parallel CP-LNS); single-skill additive 50/5/50 (against CP-LNS, Parallel CP-LNS, CP-SAT full, Greedy restarts); multi-skill additive 150/10/500 (against CP-LNS, Parallel CP-LNS, CP-SAT full, Greedy restarts); multi-skill additive 150/5/500 (against CP-LNS, Parallel CP-LNS, CP-SAT full, Greedy restarts)}
\newcommand{\CtwoOnetests}{46}
\newcommand{\CtwoOnesig}{31}
\newcommand{\CtwoOnemaxp}{1}
\newcommand{\CtwoOnewon}{361}
\newcommand{\CtwoOnepairs}{410}
\newcommand{\CtwoHalfcpsatlo}{0.953}
\newcommand{\CtwoHalfcpsathi}{1.021}
\newcommand{\CtwoHalfpcpsatlo}{0.963}
\newcommand{\CtwoHalfpcpsathi}{1.011}
\newcommand{\CtwoHalfcpfulllo}{0.940}
\newcommand{\CtwoHalfcpfullhi}{1.007}
\newcommand{\CtwoHalfctaslo}{0.136}
\newcommand{\CtwoHalfctashi}{0.568}
\newcommand{\CtwoHalfconstructlo}{0.954}
\newcommand{\CtwoHalfconstructhi}{1.019}
\newcommand{\CtwoHalfrllo}{0.765}
\newcommand{\CtwoHalfrlhi}{0.852}
\newcommand{\CtwoHalfclasslo}{-2.1}
\newcommand{\CtwoHalfclasshi}{6.0}
\newcommand{\CtwoHalfrlgainlo}{14.8}
\newcommand{\CtwoHalfrlgainhi}{23.5}
\newcommand{\CtwoHalfheld}{1}
\newcommand{\CtwoHalfverdict}{is not met}
\newcommand{\CtwoHalfsettings}{8}
\newcommand{\CtwoHalfheldlist}{multi-skill additive 50/5/200}
\newcommand{\CtwoHalffaillist}{single-skill binary 25/5/50 (against Parallel CP-LNS); single-skill binary 50/5/50 (against CP-LNS, Parallel CP-LNS, CP-SAT full); single-skill additive 50/5/50 (against CP-LNS, Parallel CP-LNS, CP-SAT full, Greedy restarts); multi-skill additive 25/5/50 (against CP-LNS, Greedy restarts); multi-skill additive 50/5/50 (against Parallel CP-LNS, Greedy restarts); multi-skill additive 150/10/500 (against CP-LNS, Parallel CP-LNS, CP-SAT full, Greedy restarts); multi-skill additive 150/5/500 (against CP-LNS, Parallel CP-LNS, CP-SAT full, Greedy restarts)}
\newcommand{\CtwoHalftests}{46}
\newcommand{\CtwoHalfsig}{26}
\newcommand{\CtwoHalfmaxp}{1}
\newcommand{\CtwoHalfwon}{348}
\newcommand{\CtwoHalfpairs}{410}
\newcommand{\Workerslo}{0.973}
\newcommand{\Workershi}{0.996}
\newcommand{\Workersgainlo}{0.4}
\newcommand{\Workersgainhi}{2.7}
\newcommand{\Cpualnslo}{0.99}
\newcommand{\Cpualnshi}{1.00}
\newcommand{\Gapalnslo}{0.0}
\newcommand{\Gapalnshi}{1.5}
\newcommand{\Cpucpsatlo}{0.28}
\newcommand{\Cpucpsathi}{0.79}
\newcommand{\Gapcpsatlo}{3.6}
\newcommand{\Gapcpsathi}{13.3}
\newcommand{\Cpupcpsatlo}{0.99}
\newcommand{\Cpupcpsathi}{1.00}
\newcommand{\Gappcpsatlo}{5.5}
\newcommand{\Gappcpsathi}{13.5}
\newcommand{\Cpucpfulllo}{0.43}
\newcommand{\Cpucpfullhi}{0.90}
\newcommand{\Gapcpfulllo}{7.8}
\newcommand{\Gapcpfullhi}{15.8}
\newcommand{\Cpuctaslo}{0.43}
\newcommand{\Cpuctashi}{0.97}
\newcommand{\Gapctaslo}{87.3}
\newcommand{\Gapctashi}{881.6}
\newcommand{\Cpuconstructlo}{1.00}
\newcommand{\Cpuconstructhi}{1.00}
\newcommand{\Gapconstructlo}{8.1}
\newcommand{\Gapconstructhi}{14.6}
\newcommand{\Cpurllo}{0.94}
\newcommand{\Cpurlhi}{0.99}
\newcommand{\Gaprllo}{27.2}
\newcommand{\Gaprlhi}{47.1}
\newcommand{\OneCoreBoneAhead}{8}
\newcommand{\RestartsRLlo}{13}
\newcommand{\RestartsRLhi}{23}
\newcommand{\CtasSolved}{12}
\newcommand{\CtasRuns}{60}
\newcommand{\LateBone}{0}
\newcommand{\RunsBone}{480}
\newcommand{\RLNlo}{3}
\newcommand{\RLNhi}{24}
\newcommand{\ItsFiftylo}{7{,}700}
\newcommand{\ItsFiftyhi}{12{,}800}
\newcommand{\ItsFivehundred}{4{,}300}
\newcommand{\RefValFreelo}{2.8}
\newcommand{\RefValFreehi}{12.0}
\newcommand{\RefValFreematched}{0}
\newcommand{\RefValInstances}{70}
\newcommand{\RefValBoundlo}{12}
\newcommand{\RefValBoundhi}{58}
\newcommand{\RefTestFreelo}{--}
\newcommand{\RefTestFreehi}{--}
\newcommand{\RefTestFreematched}{0}
\newcommand{\RefTestInstances}{0}
\newcommand{\RefTestBoundlo}{--}
\newcommand{\RefTestBoundhi}{--}
\newcommand{\RLsamplegainlo}{--}
\newcommand{\RLsamplegainhi}{--}

% generated by paper/tools/sensitivity.py --split val
\newcommand{\LateRL}{0}
\newcommand{\LateALNS}{0}
\newcommand{\LateOther}{0}
\newcommand{\Senscpsatgainlo}{3.2}
\newcommand{\Senscpsatgainhi}{11.5}
\newcommand{\Senspcpsatgainlo}{4.9}
\newcommand{\Senspcpsatgainhi}{11.1}
\newcommand{\Senscpfullgainlo}{7.0}
\newcommand{\Senscpfullgainhi}{13.3}
\newcommand{\Sensctasgainlo}{46.2}
\newcommand{\Sensctasgainhi}{87.3}
\newcommand{\Sensconstructgainlo}{6.6}
\newcommand{\Sensconstructgainhi}{11.5}
\newcommand{\Sensrlgainlo}{20.1}
\newcommand{\Sensrlgainhi}{32.0}

\renewcommand{\thetable}{S\arabic{table}}
\renewcommand{\thefigure}{S\arabic{figure}}
\renewcommand{\thesection}{S\arabic{section}}

\title{Supplementary Material\\[0.3ex]\large Search Beats Learning at Equal Compute: A Pre-Registered Study of
Heterogeneous Multi-Robot Coalition Formation and Scheduling}
\author{Anonymous submission}
\date{}

\begin{document}
\maketitle

\noindent This document accompanies the paper; section, table and equation numbers without the prefix S refer to the
paper. The archive also contains the anonymised source code of the solver, the competitors and the harness, the
pre-registration, every result row of the test campaign and the scripts that produce every table and figure.

\section{Contents of the archive and how to reproduce}
\label{sec:s-archive}

\begin{itemize}
\item \texttt{cbba\_sota/}: the exact core (instances, plans, forward pass, simulator replay), ALNS (numba kernels and
  driver), the CP-SAT baselines, the CTAS-D port and the wrapper around the released RL policy.
\item \texttt{scripts/}: instance generation and checks, the timed campaign harness (\texttt{anytime.py}), the best-known
  plans (\texttt{bks.py}) and the CTAS-D checks.
\item \texttt{docs/prereg-phase1.md}: the pre-registration as frozen (Section~\ref{sec:s-prereg}).
\item \texttt{runs/anytime\_test/}: one JSON row per run of the test campaign: instance fingerprint, method, cores,
  budget, the plan as routes, its simulator makespan, wall and CPU time, the anytime trace, the pinned CPUs, the host
  load and the container's throttling counters.
\item \texttt{docs/results/}: the reports and CSV files of every campaign (test, val, dev), and \texttt{paper/tools/}:
  the scripts that turn them into the tables and figures.
\end{itemize}
The benchmark itself (generator, simulator and released policy) is the public code of \citet{dai2025heterogeneous}; our
instances are regenerated from their seeds (\texttt{scripts/gen\_instances.py --check} verifies all 1490 files). The
test campaign runs from a clean checkout of the frozen code with
\begin{quote}\small\ttfamily
anytime.py run --split test --grid test --n 50 --n-large 50\\
\hspace*{1em}--lanes 6 --out runs/anytime\_test
\end{quote}
and \texttt{docs/results/test/analyze.sh} produces its reports.

\section{Exactness of the evaluator}
\label{sec:s-exact}

The benchmark's simulator executes pre-set routes event by event and lets an agent join a task only while the task is
not yet covered. For a plan in our representation (a global order and minimal covers), the forward pass of
Section~\ref{M-sec:repr} performs the same floating-point operations as the simulator: travel times are computed pair by
pair with the same norm, arrivals, starts and finishes in the same order. Our test suite replays hundreds of random
minimal-cover plans on all benchmark families, up to 150 robots and 500 tasks, in the simulator and requires the same
makespans (to $10^{-9}$); the plans ALNS returns are checked against the evaluator, and the incremental insertion of
Section~\ref{M-sec:insertion} against a full forward pass. Every result in the paper is the simulator's own score of the returned plan
(RL: the score of its own rollout), so an evaluator error could only have made ALNS search the wrong objective, never
report a wrong makespan.

\section{ALNS parameters}
\label{sec:s-params}

\begin{table}[h]
\centering
\small
\caption{ALNS parameters (defaults of \texttt{ALNSConfig} at the frozen commit; set on dev).}
\label{tab:s-alns}
\begin{tabular}{lll}
\toprule
Component & Parameter & Value \\
\midrule
Objective & weight $\lambda$ of the mean finish time & 0.1 \\
Construction & share of the budget & 5\% (at least one dispatch) \\
 & noise levels of randomised dispatch & 0.01, 0.03, 0.1, 0.3 \\
Destroy & tasks removed $q$ & uniform in $\{1, \dots, \mathrm{clip}(\mathrm{round}(4\sqrt{50/m}), 2, 6)\}$ \\
 & randomisation power (related, worst wait) & 6, 3 \\
Repair & positions tried per task & all, or 64 near the task \\
 & arrival noise of the noisy repair & up to 20\% \\
 & regret-2 & up to 10 removed tasks \\
Acceptance & start temperature & $0.01\, g_0 \min(1, (50/m)^{1.5})$ \\
 & end temperature & start / 50, geometric in wall time \\
Adaptation & scores (new best, better, accepted) & 33, 9, 13 \\
 & reaction factor, segment length & 0.1, 100 iterations \\
Other & re-sort by start time after acceptance & probability 0.1 \\
 & restart from the best plan & after 3000 iterations without a new best \\
 & workers, exchange period & one per core, 1\,s \\
\bottomrule
\end{tabular}
\end{table}

\section{Competitors}
\label{sec:s-competitors}

All competitors run in the same pinned lanes as ALNS; their budgets start once the instance is loaded and include
construction and model building. Table~\ref{tab:supp-params} lists the per-setting choices made on dev.

\paragraph{CP-SAT model.} One \texttt{AddCircuit} constraint per agent over its depot and its candidate tasks (tasks it
can contribute to); a skipped task is a self-loop. Arc $j \to k$ enforces $s_k \ge s_j + d_j + \tau_{jk}$, depot arcs
bound the first start and the makespan; coverage is $\sum_i a_{ik} x_{ij} \ge r_{jk}$. Valid cuts bound the number of
members per species and task, and agents of one species are ordered by the smallest task they visit (symmetry
breaking). Times are integers on a $10^{-3}$ grid, rounded up, so every CP schedule is feasible in real time; every
plan is converted to a minimal cover and scored by the simulator.

\paragraph{CP-SAT LNS.} Each step frees a window of tasks (critical chain, spatial, temporal or random), keeps every
other coalition fixed and re-solves a small model: agents that can take a freed task get a circuit over their fixed
tasks and the freed candidates, all other agents become fixed precedence chains; the incumbent is always a complete
hint. The threaded version solves each sub-model with 8 CP-SAT workers; the parallel version runs 8 single-threaded
workers on one shared incumbent (makespan, then summed depot returns).

\paragraph{CTAS-D.} Our port rebuilds the model of the benchmark's C++/Gurobi code (planner mode
\texttt{TEAMPLANNER\_CONDET}): per species a continuous flow on the task graph with binary arc indicators, shared task
start times with big-M synchronisation on every used arc, cumulative capability coverage per required trait, and the
flow rounding and path split into per-agent routes. The benchmark authors' makespan objective is not in the public
code; we recovered it from their 50 shipped Gurobi solutions as the energy cost plus 100 times the makespan. Plugged
into our model, all 50 shipped solutions satisfy every row to print precision and reproduce the shipped objective
values, and their routes replay in the simulator to the shipped makespans. Deviations: species that cannot contribute
to a task get no variables (same optimum); the energy rows, which constrain nothing at the benchmark's energy cap, are
off; CP-SAT solves the model on a $10^{-3}$ time grid; the budget includes model building. Backend choice (on the
benchmark's released 20-task test set, not blind; 10 instances, 60\,s, 8 threads): CP-SAT solved 10 of 10 with plans
4.1\% shorter than the shipped 600-s Gurobi runs, HiGHS 10 of 10 at 2.2 times their makespan, SCIP 3 of 10. On the
500-task settings the port found no solution within $B_1$ on dev and its solver held 55 and 87\,GB of memory, so it is
counted as failing there.

\paragraph{RL(s.$N$).} The released checkpoint and the benchmark's own environment, on CPU: each of 8 single-threaded
processes samples rollouts in lockstep batches (the first batch has one sample at budgets up to 10\,s, else four;
later batches fill 95\% of the remaining time at the slowest per-sample time seen so far) and the best rollout counts.
On our regenerated test instances the policy reproduces the published RL(g.) and RL(s.10) means within half a
standard deviation on all 16 settings.

\paragraph{Restarts.} Eight independent streams (seeds 0--7) of our constructions: plain list scheduling, then
alternately a noisy dispatch and a random-order best insertion; the best plan any stream finished by the budget.

\begin{table}[h]
\centering
\small
\caption{Per-setting parameters of the competitors, chosen on the development instances: time limit of one CP-SAT call in CP-LNS and in parallel CP-LNS (s), and the number of nearest candidate tasks per arc of CP-SAT full model (all: every arc).}
\label{tab:supp-params}
\begin{tabular}{lrrrr}
\toprule
Setting & $B_1$ (s) & CP-LNS call (s) & Parallel CP-LNS call (s) & CP-SAT full arcs \\
\midrule
single-skill binary 25/5/50 & 13.45 & 2 & 0.5 & all \\
single-skill binary 50/5/50 & 19.14 & 2 & 5 & 10 \\
single-skill additive 50/5/50 & 29.80 & 2 & 0.5 & all \\
multi-skill additive 25/5/50 & 16.34 & 2 & 0.5 & all \\
multi-skill additive 50/5/50 & 23.43 & 2 & 0.5 & 10 \\
multi-skill additive 50/5/200 & 75.20 & 2 & 2 & 20 \\
multi-skill additive 150/10/500 & 560.20 & 10 & 5 & 10 \\
multi-skill additive 150/5/500 & 560.50 & 2 & 2 & 5 \\
\bottomrule
\end{tabular}
\end{table}


\section{Protocol details}
\label{sec:s-protocol}

\begin{itemize}
\item \textbf{Host.} Two Intel Xeon Platinum 8480C processors (112 physical cores) in a container with a quota of 57.6
  CPUs and 330\,GB of memory. At most six lanes of 8 CPUs, one CPU per physical core, i.e.\ 48 CPUs, so the container
  stays under its quota; the only other load was a GPU keep-alive process of about 2 CPUs. RL runs on CPU.
\item \textbf{Scheduling.} Jobs of all methods are shuffled into the same lanes; an 8-core job occupies a lane, 1-core
  jobs of one budget are bundled 8 per lane. A lane starts a unit only when at most 40\% of the container's CPU periods
  were throttled over the last second; a row whose job saw more than 40\% throttled periods is disturbed, re-run and
  left out. On this host, 1-process jobs with at most 40\% throttled periods still received 87--100\% of their core.
\item \textbf{Code.} Every timed row ran from a clean checkout of a fixed commit and records it; the harness refuses a
  modified source tree, and on the test split it also refuses code whose trees differ from the frozen ones.
\item \textbf{Late runs and failures.} Pre-registered: a run that overruns its budget by more than 10\% + 0.25\,s is
  scored by the best plan its trace reached by the budget, or as a failure if it keeps no trace; a failure (no plan,
  unfinished tasks or makespan $\ge 200$) counts as 200. Late runs of the test campaign at $B_1$ on 8 cores: RL
  \LateRL, ALNS \LateALNS, others \LateOther. Scoring them by their final plans instead (descriptive), ALNS's plans are
  \Sensrlgainlo--\Sensrlgainhi\% shorter than RL(s.$N$)'s, \Senscpsatgainlo--\Senscpsatgainhi\% than CP-SAT LNS's,
  \Senspcpsatgainlo--\Senspcpsatgainhi\% than parallel CP-SAT LNS's, \Senscpfullgainlo--\Senscpfullgainhi\% than the
  CP-SAT model's and \Sensconstructgainlo--\Sensconstructgainhi\% than the restarts'.
\item \textbf{Statistics.} Per instance the ratio of makespans at the same budget and cores; the mean ratio with a
  paired bootstrap 95\% CI (10,000 resamples, seed 0); one-sided paired $t$-tests of the mean log ratio below 0; Holm
  within each family (C1 at $B_1$: 46 tests; C1 at $2B_1$: 17; C2 at 2\,s: 46).
\item \textbf{Deviations from the pre-registration.} None in the confirmatory analysis. The late-run sensitivity above
  and every val and dev result are additional and descriptive.
\end{itemize}

\section{Full results on the test split}
\label{sec:s-full}

{\small\begin{longtable}{llrrrrrrrr}
\caption{Mean gap to the best known plan (\%) on the val split per method, number of cores and budget ($^*$: some runs late, scored at the budget by their trace or as failures, as pre-registered). CPU: CPU seconds of all processes and threads / (budget $\times$ cores) at $B_1$.}\label{tab:supp-gaps}\\
\toprule
Setting & Method & 0.5 & 1 & 2 & 5 & 10 & $B_1$ & $2B_1$ & CPU \\
\midrule
\endfirsthead
\toprule
Setting & Method & 0.5 & 1 & 2 & 5 & 10 & $B_1$ & $2B_1$ & CPU \\
\midrule
\endhead
SA-BT 25/5/50 & ALNS-1 & 8.34 & 6.91 & 5.86 & 2.87 & 3.39 & 3.96 &  & 1.00 \\
 & ALNS-8 & 7.17$^*$ & 5.89 & 3.02 & 1.72 & 0.90 & 1.53 & 1.52 & 0.99 \\
 & CP-SAT LNS-8 &  &  &  &  &  & 11.55 &  & 0.79 \\
 & Par.\ CP-SAT LNS-8 & 16.64$^*$ & 13.26 & 13.30 & 11.57 & 8.59 & 9.76 & 9.26 & 0.99 \\
 & CP-SAT model-8 & 279.75$^*$ & 15.24 & 12.97 & 12.02 & 10.88 & 11.61 & 10.41 & 0.69 \\
 & CTAS-D-8 &  &  &  &  &  & 619.76 &  & 0.95 \\
 & Restarts-1 & 15.15 & 14.04 & 13.90 & 12.76 & 12.14 & 11.99 & 11.35 & 1.00 \\
 & Restarts-8 & 13.16 & 12.80 & 12.15 & 11.26 & 10.91 & 10.50 & 10.05 & 1.00 \\
 & RL(s.$N$)-8 & 659.53$^*$ & 36.63 & 35.34 & 32.43 & 29.14 & 27.21 &  & 0.99 \\
\midrule
SA-BT 50/5/50 & ALNS-1 & 5.78 & 5.08 & 3.04 & 1.65 & 1.59 & 0.68 &  & 1.00 \\
 & ALNS-8 & 4.20$^*$ & 2.42 & 1.94 & 1.42 & 0.52 & 0.27 & 0.42 & 1.00 \\
 & CP-SAT LNS-8 &  &  &  &  &  & 3.63 &  & 0.57 \\
 & Par.\ CP-SAT LNS-8 & 18.52$^*$ & 14.23$^*$ & 9.43 & 8.04 & 7.40 & 5.49 & 5.13 & 1.00 \\
 & CP-SAT model-8 & 231.95$^*$ & 17.32 & 14.53 & 12.35 & 8.55 & 7.83 & 6.46 & 0.81 \\
 & CTAS-D-8 &  &  &  &  &  & 87.33 &  & 0.97 \\
 & Restarts-1 & 15.11 & 13.20 & 11.94 & 11.08 & 10.73 & 10.60 & 9.59 & 1.00 \\
 & Restarts-8 & 12.58 & 12.29 & 11.03 & 9.81 & 9.47 & 9.20 & 8.67 & 1.00 \\
 & RL(s.$N$)-8 & 1086.14$^*$ & 879.81$^*$ & 35.34 & 34.08 & 32.02 & 32.47 &  & 0.99 \\
\midrule
SA-AT 50/5/50 & ALNS-1 & 10.21 & 9.66 & 7.53 & 7.28 & 6.00 & 3.67 &  & 1.00 \\
 & ALNS-8 & 11.58 & 9.71 & 7.04 & 3.75 & 0.94 & 0.90 & 1.53 & 1.00 \\
 & CP-SAT LNS-8 &  &  &  &  &  & 9.28 &  & 0.70 \\
 & Par.\ CP-SAT LNS-8 & 15.82 & 14.28 & 12.99 & 11.09 & 10.62 & 8.99 & 9.03 & 1.00 \\
 & CP-SAT model-8 & 127.93$^*$ & 14.02 & 13.31 & 12.08 & 11.46 & 9.50 & 8.95 & 0.69 \\
 & CTAS-D-8 &  &  &  &  &  & 647.85 &  & 0.96 \\
 & Restarts-1 & 13.47 & 12.74 & 11.87 & 10.87 & 10.13 & 9.18 & 9.02 & 1.00 \\
 & Restarts-8 & 11.76 & 10.75 & 10.30 & 9.39 & 8.59 & 8.10 & 7.65 & 1.00 \\
 & RL(s.$N$)-8 & 654.49$^*$ & 654.49$^*$ & 286.60$^*$ & 38.75 & 35.94 & 31.21 &  & 0.99 \\
\midrule
MA-AT 25/5/50 & ALNS-1 & 9.19 & 5.95 & 5.09 & 3.58 & 3.80 & 2.47 &  & 1.00 \\
 & ALNS-8 & 7.73 & 5.53 & 4.10 & 3.03 & 2.05 & 0.80 & 1.24 & 1.00 \\
 & CP-SAT LNS-8 &  &  &  &  &  & 12.13 &  & 0.75 \\
 & Par.\ CP-SAT LNS-8 & 22.57$^*$ & 20.00 & 17.46 & 15.77 & 15.03 & 13.45 & 11.72 & 1.00 \\
 & CP-SAT model-8 & 346.27$^*$ & 19.97$^*$ & 17.89 & 16.61 & 15.55 & 13.65 & 12.52 & 0.43 \\
 & CTAS-D-8 &  &  &  &  &  & 602.53 &  & 0.81 \\
 & Restarts-1 & 16.90 & 16.01 & 15.06 & 14.68 & 14.59 & 13.89 & 13.01 & 1.00 \\
 & Restarts-8 & 14.91 & 14.56 & 13.46 & 12.35 & 12.25 & 11.69 & 11.20 & 1.00 \\
 & RL(s.$N$)-8 & 580.62$^*$ & 198.87$^*$ & 41.76 & 37.87 & 36.67 & 37.83 &  & 0.99 \\
\midrule
MA-AT 50/5/50 & ALNS-1 & 7.91 & 5.98 & 5.65 & 3.73 & 1.80 & 2.72 &  & 1.00 \\
 & ALNS-8 & 6.39 & 4.64 & 3.53 & 1.33 & 0.90 & 0.73 & 1.20 & 1.00 \\
 & CP-SAT LNS-8 &  &  &  &  &  & 13.26 &  & 0.64 \\
 & Par.\ CP-SAT LNS-8 & 23.95$^*$ & 21.87 & 19.47 & 15.98 & 12.85 & 11.06 & 9.88 & 1.00 \\
 & CP-SAT model-8 & 189.87$^*$ & 21.74$^*$ & 20.11 & 18.20 & 16.95 & 14.90 & 14.15 & 0.72 \\
 & CTAS-D-8 &  &  &  &  &  & 881.65 &  & 0.87 \\
 & Restarts-1 & 19.36 & 18.20 & 17.76 & 16.91 & 15.68 & 15.10 & 14.65 & 1.00 \\
 & Restarts-8 & 16.22 & 15.74 & 15.61 & 14.71 & 13.86 & 13.24 & 12.73 & 1.00 \\
 & RL(s.$N$)-8 & 1020.44$^*$ & 1020.44$^*$ & 48.51 & 45.79 & 43.59 & 39.35 &  & 0.99 \\
\midrule
MA-AT 50/5/200 & ALNS-1 & 9.40 & 9.12 & 6.94 & 7.43 & 5.44 & 2.94 &  & 1.00 \\
 & ALNS-8 & 10.86 & 9.37 & 6.91 & 5.56 & 4.25 & 1.36 & 0.96 & 1.00 \\
 & CP-SAT LNS-8 &  &  &  &  &  & 12.77 &  & 0.64 \\
 & Par.\ CP-SAT LNS-8 & 14.98 & 14.98$^*$ & 14.98 & 14.98 & 14.98 & 12.16 & 11.39 & 1.00 \\
 & CP-SAT model-8 & 57.01$^*$ & 95.72$^*$ & 14.98$^*$ & 14.98 & 14.98 & 14.42 & 14.06 & 0.74 \\
 & CTAS-D-8 &  &  &  &  &  & 247.85 &  & 0.43 \\
 & Restarts-1 & 15.29 & 15.16 & 15.16 & 14.80 & 14.80 & 14.65 &  & 1.00 \\
 & Restarts-8 & 14.71 & 14.71 & 14.71 & 14.71 & 14.71 & 14.65 &  & 1.00 \\
 & RL(s.$N$)-8 & 291.86$^*$ & 291.86$^*$ & 291.86$^*$ & 291.86$^*$ & 153.66$^*$ & 40.54 &  & 0.94 \\
\midrule
MA-AT 150/10/500 & ALNS-1 & 12.58 & 11.57 & 10.74 & 10.23 & 8.44 & 1.53 &  & 1.00 \\
 & ALNS-8 &  &  &  &  &  & 0.00 &  & 1.00 \\
 & CP-SAT LNS-8 &  &  &  &  &  & 12.97 &  & 0.67 \\
 & Par.\ CP-SAT LNS-8 &  &  &  &  &  & 12.35 &  & 1.00 \\
 & CP-SAT model-8 &  &  &  &  &  & 14.92 &  & 0.86 \\
 & Restarts-1 & 16.64 & 16.64 & 16.31 & 15.93 & 14.86 & 13.29 &  & 1.00 \\
 & Restarts-8 & 16.48 & 15.80 & 14.92 & 14.85 & 14.08 & 12.62 &  & 1.00 \\
 & RL(s.$N$)-8 &  &  &  &  &  & 47.08 &  & 0.98 \\
\midrule
MA-AT 150/5/500 & ALNS-1 & 9.14 & 8.95 & 8.11 & 7.58 & 6.41 & 1.80 &  & 1.00 \\
 & ALNS-8 &  &  &  &  &  & 0.07 &  & 1.00 \\
 & CP-SAT LNS-8 &  &  &  &  &  & 12.77 &  & 0.28 \\
 & Par.\ CP-SAT LNS-8 &  &  &  &  &  & 12.78 &  & 1.00 \\
 & CP-SAT model-8 &  &  &  &  &  & 15.82 &  & 0.90 \\
 & Restarts-1 & 16.15 & 16.15 & 16.15 & 15.91 & 15.88 & 13.72 &  & 1.00 \\
 & Restarts-8 & 15.88 & 15.88 & 15.82 & 15.32 & 15.18 & 12.86 &  & 1.00 \\
 & RL(s.$N$)-8 &  &  &  &  &  & 40.36 &  & 0.96 \\
\bottomrule
\end{longtable}
}

\section{Reference solutions on val}
\label{sec:s-refs}

\begin{table}[h]
\centering
\small
\caption{Val split: how far the best plan of any run without an ALNS component (competitor runs of the timed campaigns and the long CP-SAT references) is above the best known plan, mean over the instances; the 300-s references on 8 cores (parallel CP-LNS and CP-SAT full model, both started from greedy plans only), their mean gap to the best known plan; and the distance of CP-SAT's certified lower bound below the best known plan.}
\label{tab:supp-refs}
\begin{tabular}{lrrrrr}
\toprule
Setting & $n$ & best w/o ALNS & Parallel CP-LNS 300\,s & CP-SAT full 300\,s & bound below \\
\midrule
single-skill binary 25/5/50 & 10 & +7.6\% & +8.1\% & +6.1\% & 46\% \\
single-skill binary 50/5/50 & 10 & +2.8\% & +3.7\% & +5.0\% & 12\% \\
single-skill additive 50/5/50 & 10 & +6.7\% & +7.2\% & +7.0\% & 45\% \\
multi-skill additive 25/5/50 & 10 & +9.4\% & +13.3\% & +12.0\% & 58\% \\
multi-skill additive 50/5/50 & 10 & +8.1\% & +9.7\% & +13.9\% & 30\% \\
multi-skill additive 50/5/200 & 10 & +11.3\% & -- & -- & -- \\
multi-skill additive 150/10/500 & 5 & +12.0\% & -- & -- & -- \\
multi-skill additive 150/5/500 & 5 & +11.5\% & -- & -- & -- \\
\bottomrule
\end{tabular}
\end{table}


The best known plans on 50 or more tasks come mostly from ALNS runs. To check that ALNS is not only close to itself,
we ran, on the first five val instances of the five 50-task settings, parallel CP-SAT LNS and the full CP-SAT model for
300\,s on 8 cores from constructions only (10 to 20 times the compute of $B_1$), and compared the best plan of any run
without an ALNS component with the best ALNS-derived plan on every val instance (Table~\ref{tab:supp-refs}). No run
without ALNS matched it on any of the \RefValInstances{} instances. On the 20-task settings the two families meet
(Table~\ref{tab:supp-small}).

\section{Ablation of the ALNS components (dev)}
\label{sec:s-ablation}

\begin{table}[h]
\centering
\small
\caption{Ablation on dev (1 worker): mean paired makespan ratio variant / final ALNS [bootstrap 95\% CI]; above 1 means the variant is worse. v1 is the configuration of an earlier phase (own random-order construction, no re-sort, older kernels); ``no portfolio'' constructs by own insertion; ``no re-sort'' never re-sorts the order by start times; ``delay-aware covers'' adds a second cover per slot (not kept: it halves the iteration rate). Blank: not run.}
\label{tab:supp-ablation}
\begin{adjustbox}{max width=\textwidth}\begin{tabular}{llcccc}
\toprule
Budget & Setting & v1 (Phase 1) & no portfolio & no re-sort & delay-aware covers \\
\midrule
$B_1$ & single-skill, binary: 25 robots, 50 tasks & 1.005 {\scriptsize [0.998, 1.013]} & 1.000 {\scriptsize [0.990, 1.010]} & 1.007 {\scriptsize [0.998, 1.016]} & 0.998 {\scriptsize [0.988, 1.009]} \\
 & single-skill, binary: 50 robots, 50 tasks & 1.007 {\scriptsize [1.001, 1.014]} & 0.999 {\scriptsize [0.995, 1.004]} & 1.004 {\scriptsize [0.999, 1.009]} & 0.995 {\scriptsize [0.988, 1.002]} \\
 & single-skill, additive: 50 robots, 50 tasks & 1.027 {\scriptsize [1.017, 1.036]} & 1.007 {\scriptsize [0.997, 1.018]} & 1.016 {\scriptsize [1.008, 1.023]} & 0.993 {\scriptsize [0.986, 1.000]} \\
 & multi-skill, additive: 25 robots, 50 tasks & 1.006 {\scriptsize [0.998, 1.016]} & 0.996 {\scriptsize [0.987, 1.004]} & 1.009 {\scriptsize [1.001, 1.018]} & 1.002 {\scriptsize [0.995, 1.008]} \\
 & multi-skill, additive: 50 robots, 50 tasks & 1.008 {\scriptsize [0.999, 1.018]} & 0.999 {\scriptsize [0.992, 1.006]} & 1.010 {\scriptsize [1.002, 1.019]} & 0.997 {\scriptsize [0.989, 1.004]} \\
 & multi-skill, additive: 50 robots, 200 tasks & 1.106 {\scriptsize [1.090, 1.123]} & 1.034 {\scriptsize [1.022, 1.047]} & 1.029 {\scriptsize [1.022, 1.036]} & 1.006 {\scriptsize [1.002, 1.010]} \\
 & multi-skill, additive: 150 robots (10 species), 500 tasks & 1.107 {\scriptsize [1.094, 1.121]} &  &  & 1.002 {\scriptsize [0.994, 1.009]} \\
 & multi-skill, additive: 150 robots (5 species), 500 tasks & 1.121 {\scriptsize [1.092, 1.148]} &  &  & 1.002 {\scriptsize [0.998, 1.006]} \\
\midrule
2\,s & single-skill, binary: 25 robots, 50 tasks & 1.020 {\scriptsize [1.011, 1.028]} & 1.001 {\scriptsize [0.992, 1.010]} & 1.014 {\scriptsize [1.007, 1.021]} & 1.006 {\scriptsize [1.002, 1.010]} \\
 & single-skill, binary: 50 robots, 50 tasks & 1.012 {\scriptsize [1.005, 1.019]} & 0.996 {\scriptsize [0.991, 1.002]} & 1.009 {\scriptsize [1.003, 1.015]} & 0.999 {\scriptsize [0.995, 1.003]} \\
 & single-skill, additive: 50 robots, 50 tasks & 1.059 {\scriptsize [1.048, 1.070]} & 1.023 {\scriptsize [1.014, 1.032]} & 1.013 {\scriptsize [1.007, 1.020]} & 0.995 {\scriptsize [0.989, 1.000]} \\
 & multi-skill, additive: 25 robots, 50 tasks & 1.018 {\scriptsize [1.008, 1.029]} & 1.006 {\scriptsize [0.996, 1.015]} & 1.019 {\scriptsize [1.008, 1.031]} & 1.011 {\scriptsize [1.005, 1.017]} \\
 & multi-skill, additive: 50 robots, 50 tasks & 1.025 {\scriptsize [1.017, 1.033]} & 1.011 {\scriptsize [1.001, 1.020]} & 1.011 {\scriptsize [1.004, 1.020]} & 1.007 {\scriptsize [1.002, 1.013]} \\
 & multi-skill, additive: 50 robots, 200 tasks & 1.166 {\scriptsize [1.148, 1.185]} & 1.088 {\scriptsize [1.073, 1.104]} & 1.019 {\scriptsize [1.012, 1.025]} & 1.003 {\scriptsize [1.000, 1.006]} \\
 & multi-skill, additive: 150 robots (10 species), 500 tasks & 1.133 {\scriptsize [1.120, 1.149]} & 1.062 {\scriptsize [1.049, 1.077]} & 1.022 {\scriptsize [1.016, 1.029]} & 0.995 {\scriptsize [0.992, 0.998]} \\
 & multi-skill, additive: 150 robots (5 species), 500 tasks & 1.149 {\scriptsize [1.123, 1.173]} & 1.078 {\scriptsize [1.057, 1.099]} & 1.022 {\scriptsize [1.017, 1.027]} & 1.003 {\scriptsize [1.000, 1.007]} \\
\bottomrule
\end{tabular}\end{adjustbox}
\end{table}


The ablation ran on dev, the tuning split, with one worker. The earlier configuration (v1) is up to 12\% worse on the
200- and 500-task settings at $B_1$; the construction portfolio matters most at short budgets on large instances; the
start-time re-sort helps by 0.4--2.9\% on the settings where it was measured at $B_1$. Eight workers instead of one improve ALNS at $B_1$ on the test
split by \Workersgainlo--\Workersgainhi\%.

\section{The 20-task settings}
\label{sec:s-small}

\begin{table}[h]
\centering
\small
\caption{Val split, three 20-task settings (descriptive; earlier campaign on another host): mean gap to the best known plan (\%) at $B_1$ on 8 cores, and the instances that CP-SAT proved optimal in 60\,s on 8 cores.}
\label{tab:supp-small}
\begin{tabular}{lrrrrr}
\toprule
Setting & ALNS & CP-LNS & Greedy restarts & RL policy & proven optimal \\
\midrule
single-skill binary 25/5/20 & 0.00 & 0.00 & 0.74 & 17.82 & 10/10 \\
single-skill additive 25/5/20 & 0.97 & 3.72 & 3.49 & 23.62 & 0/10 \\
multi-skill additive 25/5/20 & 0.16 & 4.18 & 3.53 & 29.64 & 1/10 \\
\bottomrule
\end{tabular}
\end{table}


As pre-registered, the 20-task settings are descriptive: CP-SAT proves all val instances of SA-BT 25/5/20 optimal, and
there every method except RL reaches the optimum or comes close; on the two AT settings ALNS stays closest to the best
known plans.

\section{Dynamic extension (negative result)}
\label{sec:s-dynamic}

We also explored, on dev only, a dynamic variant of the benchmark in which tasks are revealed over time and plans are
repaired online, with go/no-go gates fixed before the runs. The gate on the planner failed: our rolling planner was only
1.5\% better than construction-only replanning at equal CPU (ratio 0.985, CI [0.981, 0.989]; the gate required a ratio
of at most 0.97), so we make no dynamic claim. A comparison of communication architectures showed large headroom (the
best fixed architecture ended about 1.8 times above full information at the communication level of the gate) but did
not separate our architecture from a simple fixed rule. These results are exploratory and not part of the confirmatory
study.

\section{Pre-registration}
\label{sec:s-prereg}

The pre-registration is \texttt{docs/prereg-phase1.md} in the archive, with host names removed for anonymity. It was
frozen before any ALNS or CP-SAT variant and the CTAS-D port ran on the test split; it names the frozen code by the git
object identifiers of its source trees (withheld here for anonymity; the archive's harness checks them), and the
harness refuses to run the test split on any other code. Its disclosures are repeated in Section~\ref{M-sec:discussion}
of the paper.

\section{Use of AI assistance}
\label{sec:s-ai}

\begin{itemize}
\item \textbf{Tool and version.} Claude Code (Anthropic), version 2.1, with Anthropic Claude models (including Claude
  Opus 5.5).
\item \textbf{Roles.} The assistant wrote most of the code under the authors' direction, proposed and critiqued parts
  of the experimental design, the analysis plan and the pre-registration, ran the campaigns through the harness,
  searched and verified the bibliography against publisher and index records, and polished the text. The authors set
  the research question, made and approved the design decisions, froze the pre-registration and are responsible for
  the content.
\item \textbf{Prompts.} The prompts that shaped the experimental design are listed in \texttt{ai\_prompts.md} in the
  archive.
\end{itemize}

\bibliographystyle{ACM-Reference-Format}
\bibliography{refs}

\end{document}
